\section{Покрытие и различимость: конструктивные свидетели Rule~110}
\label{sec:observer-witnesses}

Рассмотрим два вида вопросов об одной динамике. В первом наблюдатель получает
журнал нескольких узлов и предсказывает будущую клетку. Во втором получатель
восстанавливает собственный прошлый бит $Y$ после того, как его прежнее
значение перестало быть различимым по локальному журналу. В обоих случаях
проверяется условие постоянства цели на слоях наблюдения из
раздела~\ref{sec:observer-theory}.

\subsection{Предсказуемость по журналу и предел геометрического покрытия}
\label{sec:observer-coverage}

Для начальной конфигурации $b$ обозначим её значение в позиции $s$ в момент
$t$ через $b_t(s)$. Пусть $S$ --- конечный список фиксированных позиций, а
$\tau$ --- последний момент чтения. Журнал наблюдения есть массив
\[
 H_{S,\tau}(b)=\bigl(b_k(i)\bigr)_{0\le k\le\tau,\,i\in S}.
\]
В исходнике \nolinkurl{observer/Rule110Observer.lean} предикат
\texttt{CanPredict} требует, чтобы декодер по этому журналу давал $b_t(s)$
на всех допустимых конфигурациях. Теорема \texttt{canPredict\_iff}
отождествляет это требование с постоянством $b_t(s)$ на слоях $H_{S,\tau}$.
Определение допускает произвольное ограничение начальных конфигураций.

Обозначим сдвинутую область точной зависимости через $I_t(s)=s+I_t$.
Используем установленную для $D_t=I_t$ теорему
\texttt{Rails.essential\_iff}; инвариантность динамики относительно сдвига
переносит её из начала координат в $s$. Предикат \texttt{CoversTarget} означает
\begin{equation}
 I_t(s)\ \subseteq\ \bigcup_{i\in S}\bigcup_{k=0}^{\tau} I_k(i).
 \label{eq:observer-cover}
\end{equation}
Для класса всех начальных конфигураций это необходимое условие. Если $j$
из левой части отсутствует справа, существенность $j$ даёт две конфигурации,
которые различаются только в $j$ и имеют разные $b_t(s)$. Все чтения в их
журналах совпадают, поскольку $j$ лежит вне областей зависимости этих чтений.
Такая пара нарушает критерий предсказуемости. Ограничение на класс начальных
состояний существенно: на одноэлементном классе любой будущий бит можно
предсказывать константой.

\paragraph{Контрпример к достаточности.}
Возьмём $S=(-1,1)$, $\tau=1$, $s=0$, $t=2$. Уже конусы последних чтений
равны $[-2,0]$ и $[0,2]$ и покрывают $I_2(0)=[-2,2]$.
В первой исходной конфигурации единственная единица стоит в позиции $-1$;
во второй единицы стоят в $-1$ и $0$. Вне указанных позиций в обоих случаях
нули. После одного такта множества единиц равны соответственно
$\{-2,-1\}$ и $\{-2,-1,0\}$. На датчиках в обоих мирах наблюдается один
журнал $(10,10)$. На втором такте в позиции $0$ локальные тройки $100$ и
$110$ дают соответственно $0$ и $1$. Значит, покрытие выполнено и
предсказание по полному журналу невозможно. Точные конструкции находятся в
\texttt{separated\_cells\_counterexample} и \texttt{coverage\_does\_not\_suffice}.

Эффект сохраняется и для соседних датчиков. При $S=(0,1)$, $\tau=2$ конусы
$[-2,2]$ и $[-1,3]$ покрывают $I_3(0)=[-2,3]$. Исходные множества единиц
$\{2\}$ и $\{-1,2\}$ дают один журнал $(00,01,11)$ и противоположные
значения $1$ и $0$ в позиции $0$ на такте $3$. Это
\texttt{adjacent\_cells\_counterexample}. Напротив, есть достаточное условие
с прямой интерпретацией: если на последней наблюдаемой строке прочитаны
значения \emph{всех клеток} $I_{t-\tau}(s)$ при $t\ge\tau$, то их обычная
эволюция однозначно задаёт цель. Именно потерю различимости при сжатии
начальных данных в несколько значений клеток допускает
\eqref{eq:observer-cover}.

\subsection{Четыре успешные истории и полный 11-битный контроль}

На целочисленной прямой поместим 11-битное слово в позиции $0,\ldots,10$
слева направо. Фон вне слова нулевой. Перед первым обновлением запишем
единицу в позицию $11$. Целевой бит прошлого равен $Y=b_0(5)$, показание
датчика $Z=b_2(10)$, журнал получателя $H=(b_k(0))_{k=0}^3$, а его последний
бит $K=b_3(0)$. После наблюдения на третьем такте записываем $Z$ в $-1$ и
считываем штатное обновление этой клетки на четвёртом такте.

\begin{center}
\begin{tabular}{lccccc}
Исходное слово & $H$ & $K$ & $Y$ & $Z$ & Выход на такте 4 \\
\hline
\texttt{10111010001} & 1100 & 0 & 0 & 0 & 0 \\
\texttt{10111110111} & 1100 & 0 & 1 & 1 & 1 \\
\texttt{11000010111} & 1101 & 1 & 0 & 1 & 0 \\
\texttt{11001100001} & 1101 & 1 & 1 & 0 & 1
\end{tabular}
\end{center}

На этих четырёх историях $Y=K\oplus Z$. В каждой паре с фиксированным $K$
до записи совпадает даже вся основная конфигурация на такте 3.
Левый сосед приёмной клетки равен единице, а правый равен $K$, поэтому
введённый бит превращается в правильный ответ по тождеству
\begin{equation}
 f_{110}(1,z,k)=z\oplus k.
 \label{eq:observer-xor-gate}
\end{equation}
Инверсия доставленного бита меняет ответ во всех четырёх случаях. Но сами
два подготовленных слова внутри пары различаются также вне позиции $5$.

На полном равномерном классе из $2^{11}$ слов сравним каждое слово с его
копией после переворота только $b_0(5)$. Датчик в позиции $10$ на такте 2
зависит от начальных позиций $8,\ldots,12$, а журнал клетки $0$ до такта 3
целиком определяется позицией $0$ и её трёхшаговой окрестностью.
Поэтому в каждой из 1024 пар пара наблюдений $(H,Z)$ одинакова при двух
значениях $Y$. Любой детерминированный бинарный ответ по $(H,Z)$ верен
ровно на 1024/2048 историях. Перебор \nolinkurl{observer/domain.py}
подтверждает 0/1024 изменений показания и предел 1024/2048 для декодера.
Так различаются два утверждения: причинный эффект \emph{введённого} бита
на приёмную клетку и причинный эффект \emph{прежнего} бита на датчик.

\subsection{Изолированный переворот и возврат на такте 5}
\label{sec:observer-return}

Теперь поместим указанные в таблице слова в позиции от $-6$ до $13$
включительно, сохранив нулевой фон за пределами этого интервала. В этой
конструкции наблюдатель пассивно считывает $Z=b_2(3)$. До возврата
$Y=b_0(5)$ и $K=b_4(0)$. Для каждого фиксированного $K$ два слова
различаются только в позиции $5$.

\begin{center}
\begin{tabular}{cc l cc}
$K$ & $Y$ & Слово на $[-6,13]$ & $Z$ & Выход на такте 5 \\
\hline
0 & 0 & \texttt{00000001100001000000} & 0 & 0 \\
0 & 1 & \texttt{00000001100101000000} & 1 & 1 \\
1 & 0 & \texttt{00010111010001000010} & 1 & 0 \\
1 & 1 & \texttt{00010111010101000010} & 0 & 1
\end{tabular}
\end{center}

Изолированный переворот $Y$ меняет $Z$ при обоих $K$. При $K=0$ датчик
возвращает $Y$, при $K=1$ --- его дополнение. На такте 4 в каждой паре
совпадают основные конфигурации на всей прямой. Совпадает и журнал приёмного
окна $(-2,-1,0)$ на всех тактах от 0 до 4 перед записью. В этот момент
$b_4(-2)=1$. Однократная запись $Z$ в $-1$ и одно обновление Rule~110
дают $b_5(-1)=Z\oplus K=Y$ по~\eqref{eq:observer-xor-gate}. Если записать
$1-Z$, получим $1-Y$ во всех четырёх мирах. До записи основной слой утратил
различие, но внешний однобитный регистр его сохранил. Расчёт на бесконечной
прямой реализован в \nolinkurl{observer/causal_line.py}.

\paragraph{Область минимальности.}
Такт 4 минимален для полного слияния в следующем классе на прямой: внутри
каждой пары переворачивается одна исходная клетка, снимается одно сырое
показание заранее фиксированной клетки в заранее фиксированное время, а
оба бинарных входа $K,Z$ совместно необходимы для точного восстановления
на квадрате из четырёх историй. При слиянии к такту 3 всё потенциальное
различие на этом такте лежит в интервале радиуса 3 вокруг переворота и
определяется исходным интервалом радиуса 6. Кроме самой изменяемой клетки
это 12 битов окружения. Среди $2^{12}=4096$ окружений полное слияние к
третьему такту происходит для восьми. При любом одном фиксированном
чувствительном датчике и моменте считывания до слияния ориентация его
реакции на переворот одинакова для всех этих окружений, в которых он
чувствителен. Для точного ответа датчик должен быть чувствителен в обоих
контекстах. Тогда одинаковая ориентация сделала бы $Y$ функцией одного $Z$.
Это противоречит совместной необходимости двух входов. Четыре слова из
таблицы достигают границы слияния и дают выход следующим обновлением.
Перебор окружений выполняет \nolinkurl{observer/causal_crosscheck.py}.

\subsection{Доставка с конечной скоростью и автономное исполнение}

В предыдущем опыте внешнее устройство записывало уже сохранённый бит
непосредственно в адрес приёма. Для тех же четырёх исходных слов можно
явно ограничить скорость канала. На тактах $2,3,4,5,6,7$ изолированный бит
занимает соответственно позиции $3,3,3,2,1,0$. Расписание одинаково во всех
мирах, перемещение за такт не превосходит одной клетки. На такте 7 соседи
клетки $0$ равны $b_7(-1)=1$, $b_7(1)=K$. После записи доставленного $Z$
штатное обновление даёт $b_8(0)=Y$. Замена $Z$ его дополнением обращает
все четыре ответа. До доставки журнал окна $(-1,0,1)$ попарно одинаков при
фиксированном $K$. Эту проверку выполняет \nolinkurl{observer/local_return.py}.

Автономная реализация использует кольцо длины 61 и один однородный,
не зависящий от времени локальный переход радиуса 1. Клетка хранит основной
бит, счётчик фазы, неизменную роль и один бит сообщения. При инициализации
роль наблюдателя назначена в позиции 3, приёмника --- в позиции 0.
Фаза управляет захватом, локальной пересылкой и подстановкой сообщения.
Внешних команд после инициализации нет. Для каждой пары миров основные
слои совпадают на тактах от 4 до 7, а полные расширенные состояния остаются
различными. Проверка \nolinkurl{observer/autonomous.py} даёт 4/4 верных
выхода на такте 8. Существование такого расширенного конечного алфавита
не закрывает вопрос R110.11 о том же полном протоколе внутри одной строго
бинарной однородной Rule~110. При полном совпадении двух состояний всей
изолированной детерминированной системы различие уже не может возникнуть
без нового различающего входа.

\paragraph{Статус формальной проверки.}
В историческом срезе отдельная версия наблюдательного модуля на Lean~4.33.1
дала девять отчётов \texttt{\#print axioms}, включая критерий предсказуемости
по слоям и контрпримеры к достаточности покрытия. В сборке монографии Lean~4.35.0 проверил канонический
исходник с импортами \texttt{Rule110Exact} и \texttt{Rule110FrontGlue}
и выдал все 11 из 11 запрошенных списков аксиом. Журналы приведены в
приложении~\ref{app:lean}. Для конечных примеров этого раздела отдельными
свидетельствами служат траектории и результаты перебора.
